% p153_1; source PDF page 153; manually verified editable table.
\begin{tikzpicture}[x=1cm,y=-1cm,inner sep=0pt,outer sep=0pt]
\path[use as bounding box] (0,0) rectangle (16.000,11.175);
\draw[line width=.4pt] (0.000,0.000) rectangle (1.800,0.855);
\node[align=center,text width=1.540cm,font={\linespread{1}\fontsize{10}{13.50}\selectfont\bfseries}] at (0.900,0.427) {角色};
\draw[line width=.4pt] (1.800,0.000) rectangle (5.100,0.855);
\node[align=center,text width=3.040cm,font={\linespread{1}\fontsize{10}{13.50}\selectfont\bfseries}] at (3.450,0.427) {主观要件};
\draw[line width=.4pt] (5.100,0.000) rectangle (16.000,0.855);
\node[align=center,text width=10.640cm,font={\linespread{1}\fontsize{10}{13.50}\selectfont\bfseries}] at (10.550,0.427) {范例};
\draw[line width=.4pt] (1.800,0.855) rectangle (5.100,11.175);
\node[align=center,text width=3.040cm,font={\linespread{1}\fontsize{10}{13.50}\selectfont}] at (3.450,6.015) {都要求故意};
\draw[line width=.4pt] (0.000,0.855) rectangle (1.800,3.060);
\node[align=center,text width=1.540cm,font={\linespread{1}\fontsize{10}{13.50}\selectfont}] at (0.900,1.958) {帮助犯};
\draw[line width=.4pt] (5.100,0.855) rectangle (16.000,3.060);
\node[align=left,text width=10.640cm,font={\linespread{1}\fontsize{10}{13.50}\selectfont}] at (10.550,1.958) {甲明知乙醉酒并要开车，仍将车借给乙。乙醉驾不慎轧死丙。乙构成交通肇事罪。甲存在监管过失，构成交通肇事罪的单独犯罪。在此，不能说“甲构成交通肇事罪的帮助犯”};
\draw[line width=.4pt] (0.000,3.060) rectangle (1.800,5.265);
\node[align=center,text width=1.540cm,font={\linespread{1}\fontsize{10}{13.50}\selectfont}] at (0.900,4.163) {教唆犯};
\draw[line width=.4pt] (5.100,3.060) rectangle (16.000,5.265);
\node[align=left,text width=10.640cm,font={\linespread{1}\fontsize{10}{13.50}\selectfont}] at (10.550,4.163) {乙开车送领导甲去机场。甲为了赶时间，教唆乙闯红灯，认为没事。乙照办，不慎撞死人。乙构成交通肇事罪。甲存在指挥过失，构成交通肇事罪的单独犯罪。在此不能说“甲构成交通肇事罪的教唆犯”};
\draw[line width=.4pt] (0.000,5.265) rectangle (1.800,7.470);
\node[align=center,text width=1.540cm,font={\linespread{1}\fontsize{10}{13.50}\selectfont}] at (0.900,6.368) {间接正犯};
\draw[line width=.4pt] (5.100,5.265) rectangle (16.000,7.470);
\node[align=left,text width=10.640cm,font={\linespread{1}\fontsize{10}{13.50}\selectfont}] at (10.550,6.368) {医生欺骗护士给病人打了毒针，导致病人死亡。护士构成医疗事故罪，但不能说医生构成医疗事故罪的间接正犯。医生构成故意杀人罪的间接正犯};
\draw[line width=.4pt] (0.000,7.470) rectangle (1.800,11.175);
\node[align=center,text width=1.540cm,font={\linespread{1}\fontsize{10}{13.50}\selectfont}] at (0.900,9.323) {共同正犯};
\draw[line width=.4pt] (5.100,7.470) rectangle (16.000,11.175);
\node[align=left,text width=10.640cm,font={\linespread{1}\fontsize{10}{13.50}\selectfont}] at (10.550,9.323) {甲、乙有意思联络，误将丙当作野猪打死，无法查明谁打的致命一枪。（1）根据要求故意且无罪说（少数说），甲乙不构成共同正犯，甲乙都无罪。（2）根据要求故意且有罪说（多数说），甲乙不构成共同正犯，甲乙各自构成过失致人死亡罪};
\end{tikzpicture}
